% Generated draft: requires mathematical and source review.
\documentclass[11pt]{article}
\usepackage[margin=25mm]{geometry}
\usepackage{fontspec}
\setmainfont{Latin Modern Roman}
\usepackage{amsmath,amssymb,mathtools,unicode-math}
\setmathfont{Latin Modern Math}
\usepackage{enumitem,needspace,xurl,hyperref}
\setlength{\parindent}{0pt}
\setlength{\parskip}{5pt}
\setlength{\emergencystretch}{3em}
\newcommand{\N}{\mathbb{N}}
\newcommand{\Z}{\mathbb{Z}}
\newcommand{\Q}{\mathbb{Q}}
\newcommand{\R}{\mathbb{R}}
\newcommand{\C}{\mathbb{C}}
\newcommand{\F}{\mathbb{F}}
\newcommand{\A}{\mathbb{A}}
\newcommand{\PP}{\mathbb{P}}
\newcommand{\E}{\mathbb{E}}
\newcommand{\eps}{\varepsilon}
\newcommand{\dd}{\,\mathrm{d}}
\begin{document}
\section*{UnsolvedMath ID: 701455}
\textbf{Source problem code:} catalogue entry
\subsection{Definitions and mathematical statement}
Use the standard finite-group convention that a formation is a nonempty class of finite groups, closed under isomorphisms, homomorphic images, and finite subdirect products. Here a subdirect product of groups \(H_1,\ldots,H_n\) is a subgroup of \(H_1\times\cdots\times H_n\) whose projection onto each factor is surjective. Write \(\mathrm{form}(G)\) for the smallest formation containing \(G\). A subformation of \(\mathrm{form}(G)\) is a formation contained in that class. These conventions interpret terminology not explicitly defined in the supplied record and require confirmation against the source.

For a prime \(p\), an elementary abelian \(p\)-group is a group isomorphic to \((\Z/p\Z)^d\) for some integer \(d\geq 0\). A minimal normal subgroup of a group \(H\) is a nonidentity normal subgroup containing no proper nonidentity subgroup normal in \(H\). Let \(A_5\) denote the alternating group on five letters.

The question attributed to V. P. Burichenko is the following. For every prime \(p\) and every finite group \(G\) admitting normal subgroups \(G_1,G_2\) with
\[
1<G_1<G_2<G,
\]
assume that \(G_1\) and \(G_2/G_1\) are elementary abelian groups for the same prime \(p\), that \(G_1\) is a minimal normal subgroup of \(G\), that \(G_2/G_1\) is a minimal normal subgroup of \(G/G_1\), and that
\[
G/G_2\cong A_5\times A_5.
\]
Must the set
\[
\bigl\{\mathcal H:\mathcal H\text{ is a formation and }\mathcal H\subseteq\mathrm{form}(G)\bigr\}
\]
be finite?
\subsection{Short English statement}
Does every finite group with two successive elementary abelian layers for the same prime, each minimal normal in the indicated quotient, and quotient \(A_5\times A_5\) generate a formation having only finitely many subformations? The supplied registry labels this problem open; that status has not been independently verified.
\subsection{Sources}
\begin{itemize}
\item Local UnsolvedMath registry, permanent ID \(701455\), exact source code \texttt{catalogue entry}, title ``Kourovka Notebook Problem 18.26.'' The supplied record attributes the problem to V. P. Burichenko, with proposed year 2014, and labels it open.
\item \emph{The Kourovka Notebook (Unsolved Problems in Group Theory)}, Problem 18.26, 18th issue (2014); the supplied background also identifies the 21st issue edition (2026), \texttt{arXiv:1401.0300v48}. The supplied collection description names E. I. Khukhro and V. D. Mazurov as editors. These bibliographic details are supplied attribution and have not been independently checked.
\end{itemize}
\end{document}
